\subsection{Visual Analysis of Feature Distribution}
\label{subsec:visual_analysis}

To further evaluate the quality of the synthetic samples generated by our proposed WGAN-GP-based framework and their consistency with the underlying distribution of real industrial fault data, we perform a high-dimensional feature visualization using the t-SNE \cite{Maaten2008} algorithm. The objective of this visual inspection is to verify whether the generated samples capture the intrinsic manifold structure of the training data and whether they effectively bridge the gaps in the feature space caused by class imbalance. In many predictive maintenance scenarios, sensor data often suffers from high dimensionality and non-linear relationships, making direct inspection difficult. t-SNE serves as an essential tool here to reduce these signals into a two-dimensional representation while preserving local neighbor relationships, thus allowing for a direct qualitative assessment of the generative performance.

Figure \ref{fig:tsne} illustrates the t-SNE projection of both real and synthetic samples for several critical fault categories. A visual inspection reveals several key characteristics. First, the synthetic samples (represented by open circles) are closely clustered around the centroids of the real sample distributions (represented by solid dots), indicating that the generator has successfully learned the local statistical properties of each fault class. There is no evidence of mode collapse, as the synthetic samples exhibit a diversity comparable to that of the real data, covering the breadth of the original clusters rather than concentrating on a single point in the latent space. This indicates that the gradient penalty (GP) has effectively stabilized the training process by constraining the discriminator's Lipschitz constant, preventing the generator from finding a single output that tricks the discriminator—a common failure mode in standard generative adversarial networks.

Moreover, the boundaries between different fault classes in the latent space are well-preserved. Even in regions where class overlap is naturally high due to similar physical fault signatures—such as slight bearing wear versus severe lubrication issues—the generated samples do not indiscriminately fill the inter-class margins, which would otherwise introduce noise into the subsequent classification process. Instead, the WGAN-GP generator, guided by the gradient penalty, respects the density gradients of the original distribution. This class-wise distribution alignment is crucial for effective data augmentation; it ensures that the decision boundaries learned by a downstream diagnosis model, such as a deep neural network or a support vector machine, are reinforced rather than blurred or confused. The synthetic samples act as a "buffer" that clarifies the separation between normal and faulty operation states, which is fundamentally why the diagnostic accuracy improves.

The high fidelity of the generated distribution compared to the real feature space confirms that the synthetic data can be treated as viable training candidates. By populating the under-represented regions of the feature space, particularly for rare fault modes that occur less than 1% of the time in the raw datasets, the generator provides the diagnostic classifier with sufficient boundary information to improve its generalization performance across the entire spectrum of failures. This overlap analysis demonstrates that the structural information of the mechanical system, as encoded in the raw sensor signals, is accurately reconstructed in the synthetic domain, lending strong support to the efficacy of the proposed generative strategy. Furthermore, the visual overlap between the real and synthetic clusters suggests that the generator has not only memorized the training points but has truly captured the underlying probability density function of the fault features.

\subsection{Ablation Study}
\label{subsec:ablation}

To rigorously assess the individual contributions of the core components within our proposed framework—specifically the WGAN-GP architecture and the Class-Aware Consistency (CAC) module—we conduct a systematic ablation study. We compare four specific configurations to isolate the performance gains: (1) a baseline GAN-based augmentation, (2) the standard WGAN without gradient penalty, (3) WGAN-GP without the CAC module, and (4) the full proposed architecture. All variants are evaluated on the same imbalanced test set across multiple metrics, including F1-score, Precision, Recall, and Geometric Mean (G-mean), which is particularly effective for evaluating performance on highly skewed datasets.

The quantitative results are summarized in Table \ref{tab:ablation}. Comparing the standard GAN with WGAN-GP, we observe a significant improvement in training stability and diagnostic accuracy. As shown in the loss curves in Figure \ref{fig:loss_curves}, the standard GAN exhibits significant oscillations and vanishing gradients, leading to suboptimal generator convergence where the loss function fails to provide a meaningful signal for improvement. In contrast, the WGAN-GP loss remains stable and correlates well with the visual quality of the samples. The Wasserstein distance provides a meaningful metric of convergence that prevents the generator from producing repetitive or low-quality features. The gradient penalty term ensures the Lipschitz constraint is satisfied, which is the absolute mathematical foundation for the stable gradients observed in our experiments.

The importance of the Class-Aware Consistency (CAC) module is further highlighted by the performance drop when it is removed. Without CAC, the generator occasionally produces "ambiguous" samples that lie too close to the decision boundaries of different classes. For instance, a generated sample intended to represent a "Gear Pitting" fault might inadvertently share features with a "Shaft Misalignment" fault, leading to a higher false-positive rate for neighboring fault categories. By introducing the CAC loss, which acts as a secondary supervised signal during training, the generator is forced to maximize the identity consistency of the generated samples. This ensures that each synthetic sample is clearly attributable to its target fault class according to a pre-trained auxiliary classifier. This leads to a marked increase in the G-mean, particularly for minor classes that are frequently confused with the "Normal" operation state.

Our results show that while WGAN-GP improves the overall quality of the data, the CAC module is responsible for nearly 40% of the accuracy gain in the minority classes. Overall, the ablation results confirm that while WGAN-GP provides the fundamental stability required for industrial signal generation, the CAC module is the key driver for achieving high classification precision in multi-class scenarios. The synergy between these components allows the framework to generate not just "realistic" signals, but "informative" ones that specifically target the weaknesses of the classifier under imbalanced conditions. The combination of structural stability from the GAN architecture and class-specific guidance from the consistency module results in a balanced dataset that significantly boosts the performance or any standard diagnostic model.

\subsection{Sensitivity Analysis and Robustness}
\label{subsec:sensitivity}

In this final part of the experimental evaluation, we investigate the robustness of the proposed framework under varying operational constraints, specifically focusing on its sensitivity to the initial training sample size and the severity of the class imbalance. An ideal diagnostic system must maintain high performance even when the available labeled data for certain faults is extremely scarce, a condition frequently encountered in the early stages of industrial deployment or for catastrophic failures that rarely occur.

First, we vary the number of real training samples available for the minority classes from 5\% to 50\% of the majority class size to test the limits of the data augmentation component. Figure \ref{fig:sensitivity} shows the change in diagnostic accuracy as a function of this imbalance ratio. The results indicate that our method maintains an F1-score above 0.85 even when the imbalance ratio is as severe as 1:20 (5%). This robustness is attributed to the framework’s ability to extract deep latent features from fewer samples and expand them into a representative synthetic set. While performance naturally scales with the amount of real data, the rate of decline in our method is significantly flatter compared to traditional Oversampling (SMOTE) or standard GAN-based methods. This suggests that the WGAN-GP is less prone to "overfitting" on the small minority set, as the Wasserstein objective encourages a smoother transition across the feature space than the linear interpolation used in SMOTE.

Furthermore, we evaluate the robustness of the Sparse Autoencoder (SAE) feature extractor used in our diagnostic pipeline against external environmental noise. Figure \ref{fig:sae_error} tracks the SAE reconstruction error variation across different noisy environments where the Signal-to-Noise Ratio (SNR) varies from -5dB to 20dB. The proposed model exhibits lower sensitivity to Gaussian noise injected into the input signals compared to traditional deep learning models. This is likely because the generative augmentation process itself acts as a form of data-centric regularization, teaching the model to ignore non-essential fluctuations and focus on the core structural components of the fault signal that are invariant to noise.

Table \ref{tab:robustness} provides the detailed results of these robustness tests across different operating speeds (800, 1200, and 1500 RPM) and loads (0, 10, and 20 Nm). The consistent performance across these varying conditions—with a total accuracy variance of less than 2.3\%—suggests that the combination of WGAN-GP for distribution matching and the SAE for robust feature extraction creates a diagnostic system that is highly adaptable to the complexities of real-world industrial environments. The ability to handle varying degrees of data scarcity, environmental noise, and varying load conditions makes this framework a reliable and practical solution for automated fault diagnosis in critical machinery where "perfect" data is rarely available. The stability of the SAE reconstruction error specifically proves that the learned features remain invariant to operational shifts, which is a key requirement for deployment in variable-speed industrial processes.
